% Part: turing-machines
% Chapter: machines-computations
% Section: variants

\documentclass[../../../include/open-logic-section]{subfiles}

\begin{document}

\olfileid{tur}{mac}{var}
\olsection{Turing机的变体}

事实上，定义图灵机有多种可能的方式，我们的定义只是其中之一。
在某些方面，我们的定义比其他定义更宽松。
我们允许任意的有限字母表，而限制更多的定义可能只允许两种纸带符号，$\TMstroke$
和~$\TMblank$。
我们允许机器同时写入符号并移动读写头，其他定义则可能只允许写入或移动其一。
我们允许只写入而不移动读写头，其他定义则省略了 $\TMstay$ 这一“指令”。
在其他方面，我们的定义则更严格。
我们假设纸带仅在单方向无限延伸，而其他定义允许纸带向左和向右都无限延伸。
事实上，甚至可以允许任意多条独立的纸带，甚至是一个无限的方格网格。
我们用转移函数来表示图灵机的指令集；其他定义则使用转移关系，使得机器在任何给定情况下都可能有多于一条可用的指令。

对定义的最后这一放宽尤其有趣。
在我们的定义中，当机器处于状态~$q$ 且读到符号~$\sigma$ 时，$\delta(q, \sigma)$ 确定了新的符号、状态及读写头位置。
但如果我们允许指令集是当前状态-符号对 $\tuple{q, \sigma}$ 与新的状态-符号-方向三元组 $\tuple{q', \sigma', D}$ 之间的关系，那么图灵机的动作就可能不是唯一确定的——指令关系可能同时包含 $\tuple{q,
  \sigma, q', \sigma', D}$ 和 $\tuple{q, \sigma, q'', \sigma'', D'}$。
这种情况下，我们得到的就是一台\emph{非确定性的}图灵机。
它们在计算复杂性理论中扮演着重要角色。

关于图灵机何时停机，也有不同的约定：
我们规定当转移函数未定义时它停机，而其他定义则要求机器处于一个特殊指定的停机状态。
我们已在 \olref[hal]{sec} 中解释了为何要求一个指定的停机状态并不会影响图灵机能计算什么。
由于我们的图灵机的纸带是单向无限的，因此在某些情况下图灵机无法正确执行指令：例如它读到最左边的方格且应向左移动时。
根据我们的定义，它只是留在原处而不是“掉下纸带”，但我们本可以将其定义为在此时停机。
这种定义也是等价的：我们可以通过删除所有 $\delta(q,\TMendtape) =
\tuple{q',\sigma,\TMleft}$ 的转移，来模拟一台在试图从第~$0$ 号方格左移时就停机的图灵机的行为——这样一来，机器在遇到~$\TMendtape$ 时就不会试图左移，而是停机。\footnote{这并不\emph{完全}可行，因为没有什么能阻止我们在除第~$0$ 号方格以外的其他方格上写入和读取 $\TMendtape$（参见 \olref[una]{ex:mover}）。我们可以通过添加第二个符号~$\TMendtape'$ 来用于此目的，从而绕过这个问题。}

表示数字（进而表示图灵机计算的输入-输出函数）也有不同的方式：我们使用一元表示，但也可以使用二进制表示。
除了需要 $\TMblank$ 和~$\TMendtape$ 之外，这还需要另外两个符号。

现在有一个有趣的事实：所有这些变体都不影响哪些函数是图灵可计算的。
\emph{如果一个函数根据某一定义是图灵可计算的，那么根据所有定义，它都是图灵可计算的。}

我们不会深入探讨验证此点的细节。这里仅举一例：允许双向无限的纸带或多条纸带，并不会增加额外的计算能力。
原因大致是，一台单带单向无限纸带的图灵机可以模拟多条纸带或双向无限的纸带。
例如，利用额外的状态和指令，我们可以将一个为多纸带或双向无限纸带机器编写的程序“翻译”成单条单向无限纸带机器的程序。
翻译后的机器可以用偶数方格对应纸带~$1$ 的方格（或双向无限纸带的“正数”方格），用奇数方格对应纸带~$2$ 的方格（或“负数”方格）。

\end{document}